\par\smallskip
\noindent\emph{Acknowledgment.}
We thank Sahibzada Allahyar for making his reconstruction available.
It helped us follow the passage from the theta lifts to the prime-sum
estimate \cite{allahyar}.

\begingroup\small
\begin{thebibliography}{99}
\phantomsection
\addcontentsline{toc}{section}{\refname}
\markboth{\refname}{\refname}
\setlength{\itemsep}{1pt}

\bibitem{announcement}
ai agent.
\href{https://alpo.ge/remark.pdf}{Public announcement on the Riemann
hypothesis}, 20 November 2026.

\bibitem{source}
ai agent and Claude.
Formal proof accompanying the announcement, 20 November 2026. Lean~4.33.1.
\href{https://github.com/TomasOrtega/RH/tree/a6ac91a412dcbaab31801355ec5707024a71bf65/extracted/lean}{Public source}.

\bibitem{allahyar}
Sahibzada Allahyar.
\href{https://github.com/sahibzada-allahyar/math/blob/rh-exposition-2026-11-20-r2/preprint/alpoge-riemann/reconstruction.pdf}{\emph{From theta lifts to zeta zeros: A reconstruction of ai agent's announced argument}}.
Expository manuscript, revision~2, 20 November 2026.

\bibitem{f1thorner-zaman}
Jesse Thorner and Asif Zaman.
Refinements to the prime number theorem for arithmetic progressions.
\href{https://doi.org/10.1007/s00209-023-03414-3}{\emph{Mathematische
Zeitschrift} 306 (2024), no.~3, article~54}.
\href{https://arxiv.org/pdf/2108.10878v2\#page=1}{arXiv v2, Theorem~1.1},
\href{https://arxiv.org/pdf/2108.10878v2\#page=2}{(1.2) and (1.7)},
\href{https://arxiv.org/pdf/2108.10878v2\#page=3}{Corollary~1.4 and Remark~1.5}.

\bibitem{f1thorner-zaman-density}
Jesse Thorner and Asif Zaman.
An explicit version of Bombieri's log-free density estimate and
S\'ark\H{o}zy's theorem for shifted primes.
\href{https://doi.org/10.1515/forum-2023-0091}
{\emph{Forum Mathematicum} 36 (2024), 1059--1080}.
\href{https://arxiv.org/pdf/2208.11123v4\#page=2}
{arXiv v4, Theorem~1.2} and
\href{https://arxiv.org/pdf/2208.11123v4\#page=1}{(1.1)--(1.2)}.

\bibitem{f1heath-brown}
D. R. Heath-Brown.
A mean value estimate for real character sums.
\href{https://doi.org/10.4064/aa-72-3-235-275}
{\emph{Acta Arithmetica} 72 (1995), 235--275}.
\href{https://www.impan.pl/en/publishing-house/journals-and-series/acta-arithmetica/all/72/3/108575/a-mean-value-estimate-for-real-character-sums}
{Corollary~1, p.~237, and Theorem~3, p.~240}.

\bibitem{f1ihara-murty-shimura}
Yasutaka Ihara, V. Kumar Murty and Mahoro Shimura.
On the logarithmic derivatives of Dirichlet \(L\)-functions at \(s=1\).
\href{https://doi.org/10.4064/aa137-3-6}
{\emph{Acta Arithmetica} 137 (2009), 253--276}.
\href{https://www.impan.pl/shop/en/publication/transaction/download/product/82940\#page=4}
{\S2.1, p.~256}, and
\href{https://www.impan.pl/shop/en/publication/transaction/download/product/82940\#page=5}
{Theorem~1, (2.2.1), p.~257}.

\bibitem{f1bennett-zeros}
Michael A. Bennett, Greg Martin, Kevin O'Bryant and Andrew Rechnitzer.
Counting zeros of Dirichlet \(L\)-functions.
\href{https://doi.org/10.1090/mcom/3599}
{\emph{Mathematics of Computation} 90 (2021), 1455--1482}.
\href{https://personal.math.ubc.ca/~bennett/BeMaObRe-MathComp-2021.pdf\#page=2}
{Theorem~1.1 and Corollary~1.2, p.~1456}.

\bibitem{f1montgomery-vaughan-moments}
Hugh L. Montgomery and Robert C. Vaughan.
Mean values of character sums.
\href{https://doi.org/10.4153/CJM-1979-053-2}
{\emph{Canadian Journal of Mathematics} 31 (1979), 476--487}.
\href{https://www.cambridge.org/core/services/aop-cambridge-core/content/view/9372DCDEE648FF605F66A10314977F5F/S0008414X00012116a.pdf/div-class-title-mean-values-of-character-sums-div.pdf\#page=5}
{Lemma~9, (7), p.~480}.

\bibitem{f1page}
A. Page.
On the number of primes in an arithmetic progression.
\href{https://londmathsoc.onlinelibrary.wiley.com/doi/10.1112/plms/s2-39.1.116}
{\emph{Proceedings of the London Mathematical Society} (2) 39 (1935),
116--141}.
For the Landau--Page theorem in the form used here, see
\href{https://arxiv.org/pdf/1412.5029v3\#page=32}
{Ford, Green, Konyagin, Maynard and Tao, \emph{Long gaps between
primes}, arXiv v3, Lemma~7.1}.

\bibitem{f1lamzouri-languasco}
Youness Lamzouri and Alessandro Languasco.
Small values of \(\lvert L'/L(1,\chi)\rvert\).
\href{https://doi.org/10.1080/10586458.2021.1927255}
{\emph{Experimental Mathematics} 32 (2023), 362--377}.
\href{https://arxiv.org/pdf/2005.10714v1\#page=8}
{arXiv v1, Proposition~3.2, p.~8}.

\bibitem{f1milne-cft}
J. S. Milne.
\emph{Class Field Theory}, version~4.03, 2020.
Ch.~V, \href{https://www.jmilne.org/math/CourseNotes/CFT.pdf\#page=165}
{Examples~3.1--3.2, pp.~156--157},
and \href{https://www.jmilne.org/math/CourseNotes/CFT.pdf\#page=168}
{Example~3.11, p.~159}.

\bibitem{f1milne-adt}
J. S. Milne.
\emph{Arithmetic Duality Theorems}, 2nd ed., BookSurge, 2006.
\href{https://jmilne.org/math/Books/ADTnot.pdf\#page=65}
{Theorem~I.4.10(a), p.~57},
\href{https://jmilne.org/math/Books/ADTnot.pdf\#page=70}
{Corollary~I.4.15, p.~62}, and
\href{https://jmilne.org/math/Books/ADTnot.pdf\#page=178}
{Proposition~II.2.9, pp.~170--171}.

\bibitem{f1weibel-kbook}
Charles A. Weibel.
\emph{The \(K\)-book: An Introduction to Algebraic \(K\)-theory}.
Graduate Studies in Mathematics 145, American Mathematical Society, 2013.
\href{https://sites.math.rutgers.edu/~weibel/Kbook/Kbook.VI.pdf\#page=47}
{Theorem~VI.8.2} and
\href{https://sites.math.rutgers.edu/~weibel/Kbook/Kbook.VI.pdf\#page=48}
{Theorem~VI.8.4}.

\bibitem{f1lim-wild}
Meng Fai Lim.
On the codescent of \'etale wild kernels in \(p\)-adic Lie extensions.
\href{https://doi.org/10.1215/21562261-2024-0016}
{\emph{Kyoto Journal of Mathematics} 65 (2025), 315--347}.
\href{https://arxiv.org/pdf/2101.06695v2\#page=20}
{arXiv v2, Proposition~5.1, p.~20}.

\bibitem{f1bayer-fluckiger-emery-houriet}
Eva Bayer-Fluckiger, Vincent Emery and Julien Houriet.
Hermitian lattices and bounds in \(K\)-theory of algebraic integers.
\href{https://doi.org/10.4171/DMS/7/4}
{\emph{Documenta Mathematica}, Extra Volume Merkurjev (2015), 71--83}.
\href{https://ftp.gwdg.de/pub/misc/EMIS/journals/DMJDMV/vol-merkurjev/bayer_emery_houriet.pdf\#page=2}
{Theorem~1.1, p.~72}.

\bibitem{f1burns-flach}
David Burns and Matthias Flach.
On the equivariant Tamagawa number conjecture for Tate motives, Part II.
\href{https://doi.org/10.4171/DMS/4/4}
{\emph{Documenta Mathematica}, Extra Volume Coates (2006), 133--163}.
\href{https://ems.press/content/book-chapter-files/27466?nt=1\#page=2}
{Corollary~1.2, p.~134}.

\bibitem{f1soule-torsion}
Christophe Soul\'e.
A bound for the torsion in the \(K\)-theory of algebraic integers.
\emph{Documenta Mathematica}, Extra Volume Kato (2003), 761--788.
\href{https://ems.press/content/book-chapter-files/27335\#page=26}
{\S5.4, (26), p.~786}.

\bibitem{gqt}
Wee Teck Gan, Yannan Qiu and Shuichiro Takeda.
The regularized Siegel--Weil formula (the second term identity) and the
Rallis inner product formula.
\href{https://doi.org/10.1007/s00222-014-0509-0}{\emph{Inventiones
Mathematicae} 198 (2014), 739--831}.
arXiv v3, \href{https://arxiv.org/pdf/1207.4709v3\#page=14}{\S3.1},
\href{https://arxiv.org/pdf/1207.4709v3\#page=33}{\S7.1 and Theorem~7.1},
\href{https://arxiv.org/pdf/1207.4709\#page=34}{Theorem~7.3(i), the residue for a complementary pair},
\href{https://arxiv.org/pdf/1207.4709v3\#page=52}{\S\S11.3--11.6}, and
\href{https://arxiv.org/pdf/1207.4709v3\#page=53}{Propositions~11.1(ii) and~11.2}.

\bibitem{p2eischen-doubling}
Ellen Eischen, Michael Harris, Jianshu Li and Christopher Skinner.
\(p\)-adic \(L\)-functions for unitary groups.
\href{https://doi.org/10.1017/fmp.2020.4}
{\emph{Forum of Mathematics, Pi} 8 (2020), e9}.
\href{https://www.cambridge.org/core/services/aop-cambridge-core/content/view/E1D22AA60E10A5B5CA5F838EB48467D7/S2050508620000049a.pdf/p_adic_l_functions_for_unitary_groups.pdf\#page=13}
{Hypotheses~1.4.4(1), p.~13}, and
\href{https://www.cambridge.org/core/services/aop-cambridge-core/content/view/E1D22AA60E10A5B5CA5F838EB48467D7/S2050508620000049a.pdf/p_adic_l_functions_for_unitary_groups.pdf\#page=49}
{\S4.2.1, (53), p.~49}.

\bibitem{seesaw}
Stephen S. Kudla.
Seesaw dual reductive pairs.
In \emph{Automorphic Forms of Several Variables}, edited by
I.~Satake and Y.~Morita, Progress in Mathematics 46,
Birkh\"auser, 1984, pp.~244--268.
\href{https://wiki.epfl.ch/waldspurger/documents/1984\%20kudla.pdf\#page=3}
{Section~1, Definition~1.3 and (1.14), p.~249}.

\bibitem{ggp}
Wee Teck Gan, Benedict H. Gross and Dipendra Prasad.
Symplectic local root numbers, central critical \(L\)-values, and
restriction problems in the representation theory of classical groups.
\href{https://www.numdam.org/item/AST_2012__346__1_0/}
{\emph{Ast\'erisque} 346 (2012), 1--109}.
Sections~12 and~23 treat Bessel and Fourier--Jacobi coefficients.

\bibitem{depth-endoscopic-ggp}
Rapha\"el Beuzart-Plessis, Pierre-Henri Chaudouard and Micha{\l} Zydor.
The global Gan--Gross--Prasad conjecture for unitary groups:
the endoscopic case.
\href{https://www.numdam.org/articles/10.1007/s10240-021-00129-1/}
{\emph{Publications Math\'ematiques de l'IH\'ES} 135 (2022), 183--336}.
\href{https://numdam.org/item/10.1007/s10240-021-00129-1.pdf\#page=3}
{Theorems~1.1.5.1 and~1.1.5.3}.

\bibitem{depth-he}
Hongyu He.
On the Gan--Gross--Prasad conjecture for \(U(p,q)\).
\href{https://doi.org/10.1007/s00222-017-0720-x}
{\emph{Inventiones Mathematicae} 209 (2017), 837--884}.
\href{https://arxiv.org/pdf/1508.02032v2\#page=3}
{Theorem~0.1 and the signed interlacing rule in the author's version}.

\bibitem{depth-harris}
Michael Harris.
Testing rationality of coherent cohomology of Shimura varieties.
In J.~Cogdell et al., eds.,
\emph{Automorphic Forms and Related Geometry},
Contemporary Mathematics 614 (2014), 81--95.
\href{https://webusers.imj-prg.fr/~michael.harris/testing.pdf\#page=7}
{Proposition~2.6 in the author's November 2012 version}.

\bibitem{p2bp-chaudouard-bessel}
Rapha\"el Beuzart-Plessis and Pierre-Henri Chaudouard.
The global Gan--Gross--Prasad conjecture for unitary groups. II.
From Eisenstein series to Bessel periods.
\href{https://doi.org/10.1017/fmp.2025.8}
{\emph{Forum of Mathematics, Pi} 13 (2025), e16, 1--98}.
\href{https://www.cambridge.org/core/services/aop-cambridge-core/content/view/2E9DFB9FC7A671B5530E47E3A37D1E7E/S2050508625000083a.pdf/global_gangrossprasad_conjecture_for_unitary_groups_ii_from_eisenstein_series_to_bessel_periods.pdf\#page=9}
{\S\S1.3.5--1.3.6, Theorem~1.3.6.1 and (1.3.6.1), pp.~8--9}.

\bibitem{p2harris-adjoint}
Michael Harris.
\(L\)-functions and periods of adjoint motives.
\href{https://doi.org/10.2140/ant.2013.7.117}
{\emph{Algebra \& Number Theory} 7 (2013), 117--155}.
\href{https://msp.org/ant/2013/7-1/ant-v7-n1-p05-p.pdf\#page=22}
{Lemma~2.5.4, pp.~137--138}.

\bibitem{p2beuzart}
Rapha\"el Beuzart-Plessis.
A local trace formula for the Gan--Gross--Prasad conjecture for
unitary groups: the Archimedean case.
\emph{Ast\'erisque} 418 (2020).
arXiv v3, \href{https://arxiv.org/pdf/1506.01452v3\#page=144}{\S6.2},
\href{https://arxiv.org/pdf/1506.01452v3\#page=154}{Lemma~6.5.1(ii)},
\href{https://arxiv.org/pdf/1506.01452v3\#page=177}{\S7.1} and
\href{https://arxiv.org/pdf/1506.01452v3\#page=181}{Theorem~7.2.1}.

\bibitem{p2borade-et-al}
Neelima Borade, Jonas Franzel, Johannes Girsch, Wei Yao, Qiyao Yu
and Elad Zelingher.
On tori periods of Weil representations of unitary groups.
\href{https://doi.org/10.1007/s00029-025-01047-4}
{\emph{Selecta Mathematica} (N.S.) 31 (2025), article~49}.
Sections~3.1.1 and~4.1.1.

\bibitem{p2jiang-sun-zhu}
Dihua Jiang, Binyong Sun and Chen-Bo Zhu.
Uniqueness of Bessel models: the Archimedean case.
\href{https://doi.org/10.1007/s00039-010-0077-4}
{\emph{Geometric and Functional Analysis} 20 (2010), 690--709}.
\href{https://bpb-us-w2.wpmucdn.com/blog.nus.edu.sg/dist/3/12136/files/2019/08/Author-copy-GAFA-Bessel-models.pdf\#page=2}
{Theorem~A, pp.~690--691}.

\bibitem{bannai-kobayashi-theta}
Kenichi Bannai and Shinichi Kobayashi.
Algebraic theta functions and the \(p\)-adic interpolation of
Eisenstein--Kronecker numbers.
\href{https://doi.org/10.1215/00127094-2010-024}
{\emph{Duke Mathematical Journal} 153 (2010), 229--295}.
\href{https://arxiv.org/pdf/math/0610163v4}
{Example~1.9, Corollary~2.2 and Theorem~2.9 in the authors' version}.

\bibitem{tate-thesis}
John Tate.
Fourier analysis in number fields and Hecke's zeta-functions.
In J.\ W.\ S. Cassels and A. Fr\"ohlich, eds.,
\href{https://www.math.ucsd.edu/~csorense/teaching/math205/Cassels-Frohlich-Algebraic_Number_Theory.pdf}
{\emph{Algebraic Number Theory}}, Thompson Book Company, 1967,
Chapter~XV, pp.~305--347, especially Section~4 on the global zeta integrals.

\bibitem{blasius-hecke}
Don Blasius.
On the critical values of Hecke \(L\)-series.
\href{https://annals.math.princeton.edu/1986/124/p02}
{\emph{Annals of Mathematics} (2) 124 (1986), 23--63}.

\bibitem{lin-cm-periods}
Jie Lin.
Period relations for automorphic induction and applications, I.
\href{https://arxiv.org/abs/1511.03517v1}{arXiv:1511.03517v1 (2015)}.
Section~1.2, Proposition~1.1, equation~(1.1), and Theorem~1.1.

\bibitem{blasius-period-relations}
Don Blasius.
Period relations and critical values of \(L\)-functions.
\href{https://msp.org/pjm/1997/181-3/pjm-v181-n3-p03-s.pdf}
{\emph{Pacific Journal of Mathematics} 181 (1997), no.~3, 53--83},
Sections L.6--7 and H.6.

\bibitem{p2xue-bessel}
Hang Xue.
Bessel models for unitary groups and Schwartz homology.
\href{https://doi.org/10.1353/ajm.2026.a997141}
{\emph{American Journal of Mathematics} 148 (2026), no.~4, 1103--1150}.
Accepted manuscript, \href{https://preprint.press.jhu.edu/ajm/sites/default/files/AJM-h-xue.pdf\#page=2}
{\S\S1.1--1.2} and
\href{https://preprint.press.jhu.edu/ajm/sites/default/files/AJM-h-xue.pdf\#page=5}{Theorem~1.1}.
For the tempered case, see the same author's
\href{https://doi.org/10.1215/00127094-2022-0018}
{\emph{Bessel models for real unitary groups: The tempered case},
Duke Mathematical Journal 172 (2023), 995--1031}.

\bibitem{p2atobe-hc}
Hiraku Atobe.
On the non-vanishing of theta liftings of tempered representations of
\(U(p,q)\).
\href{https://doi.org/10.1016/j.aim.2020.106984}
{\emph{Advances in Mathematics} 363 (2020), article~106984}.
\href{https://eprints.lib.hokudai.ac.jp/repo/huscap/all/84508/Adv.\%20math\%20363_106984.pdf\#page=14}
{\S2.2 and Theorem~2.1(4), (6), (7), pp.~13--15 of the accepted manuscript},
\href{https://eprints.lib.hokudai.ac.jp/repo/huscap/all/84508/Adv.\%20math\%20363_106984.pdf\#page=43}
{Appendix~A.3, p.~42},
\href{https://eprints.lib.hokudai.ac.jp/repo/huscap/all/84508/Adv.\%20math\%20363_106984.pdf\#page=46}
{Appendix~A.4, p.~45}, and
\href{https://eprints.lib.hokudai.ac.jp/repo/huscap/all/84508/Adv.\%20math\%20363_106984.pdf\#page=49}
{Theorem~A.4, p.~48}.

\bibitem{p2ggp-examples}
Wee Teck Gan, Benedict H. Gross and Dipendra Prasad.
Restrictions of representations of classical groups: examples.
\href{https://www.numdam.org/book-part/AST_2012__346__111_0/}
{\emph{Ast\'erisque} 346 (2012), 111--170}.
\href{https://www.numdam.org/item/AST_2012__346__111_0.pdf\#page=5}
{Proposition~2.1, p.~114}.

\bibitem{p2ichino-real-theta}
Atsushi Ichino.
Theta lifting for tempered representations of real unitary groups.
\href{https://doi.org/10.1016/j.aim.2022.108188}
{\emph{Advances in Mathematics} 398 (2022), article~108188}.
arXiv v2, \href{https://arxiv.org/pdf/2008.06174v2\#page=10}
{\S4.1, p.~10}, and
\href{https://arxiv.org/pdf/2008.06174v2\#page=25}
{Proposition~7.1, Remark~7.2 and \S7.2, pp.~25--26}.

\bibitem{tr3getzhahn}
Jayce R. Getz and Heekyoung Hahn.
A general simple relative trace formula.
\href{https://doi.org/10.2140/pjm.2015.277.99}
{\emph{Pacific Journal of Mathematics} 277 (2015), 99--118}.
\href{https://msp.org/pjm/2015/277-1/pjm-v277-n1-p04-p.pdf\#page=18}
{\S4C, (4C.8) and the following unfolding, p.~114}.

\bibitem{tr3kowalski-saha-tsimerman}
Emmanuel Kowalski, Abhishek Saha and Jacob Tsimerman.
Local spectral equidistribution for Siegel modular forms and applications.
\href{https://doi.org/10.1112/S0010437X11007391}
{\emph{Compositio Mathematica} 148 (2012), 335--384}.
\href{https://arxiv.org/pdf/1010.3648v4\#page=18}
{arXiv v4, Proposition~3.3, p.~18 (published p.~355)}.

\bibitem{tr3shin-templier}
Sug Woo Shin and Nicolas Templier.
Sato--Tate theorem for families and low-lying zeros of automorphic
\(L\)-functions.
\href{https://link.springer.com/article/10.1007/s00222-015-0583-y}
{\emph{Inventiones Mathematicae} 203 (2016), 1--177}.
\href{https://link.springer.com/content/pdf/10.1007/s00222-015-0583-y.pdf\#page=111}
{Theorem~9.16, p.~111},
\href{https://link.springer.com/content/pdf/10.1007/s00222-015-0583-y.pdf\#page=112}
{Theorem~9.19, p.~112}, and
\href{https://link.springer.com/content/pdf/10.1007/s00222-015-0583-y.pdf\#page=118}
{Corollary~9.25, p.~118}.

\bibitem{tr3wei-poincare}
Zhining Wei.
Linear relations of Siegel Poincar\'e series and non-vanishing of
the central values of spinor \(L\)-functions.
\href{https://www.numdam.org/articles/10.5802/jtnb.1226/}
{\emph{Journal de Th\'eorie des Nombres de Bordeaux} 34 (2022), 755--766}.
\href{https://numdam.org/item/10.5802/jtnb.1226.pdf\#page=7}
{Lemma~2.3, p.~760}, and
\href{https://numdam.org/item/10.5802/jtnb.1226.pdf\#page=9}
{\S3.1, p.~762}.

\bibitem{tr3green}
J.~A. Green.
The characters of the finite general linear groups.
\href{https://doi.org/10.1090/S0002-9947-1955-0072878-2}
{\emph{Transactions of the American Mathematical Society} 80 (1955), 402--447}.
\href{https://www.ams.org/tran/1955-080-02/S0002-9947-1955-0072878-2/S0002-9947-1955-0072878-2.pdf\#page=38}
{Lemma~7.7 and Theorems~12--14, pp.~439--444}.

\bibitem{tr3gelfand}
S.~I. Gel'fand.
Representations of the full linear group over a finite field.
\href{https://doi.org/10.1070/SM1970v012n01ABEH000907}
{\emph{Mathematics of the USSR-Sbornik} 12 (1970), 13--39}.
\href{https://www.mathnet.ru/eng/sm3499}
{English translation, Theorem~0.5 and Proposition~1.2, pp.~17--18}.

\bibitem{tr3liu}
Yifeng Liu.
Arithmetic theta lifting and \(L\)-derivatives for unitary groups, I.
\emph{Algebra \& Number Theory} 5 (2011), 849--921.
\href{https://msp.org/ant/2011/5-7/ant-v5-n7-p01-p.pdf\#page=12}
{Theorem~2.1(1), p.~859}.

\bibitem{tr3yamana}
Shunsuke Yamana.
The Siegel--Weil formula for unitary groups.
\href{https://doi.org/10.2140/pjm.2013.264.235}
{\emph{Pacific Journal of Mathematics} 264 (2013), 235--257}.
\href{https://msp.org/pjm/2013/264-1/pjm-v264-n1-p11-p.pdf\#page=4}
{Theorem~2, p.~237}, and
\href{https://msp.org/pjm/2013/264-1/pjm-v264-n1-p11-p.pdf\#page=13}
{Lemma~3.1, p.~246}.

\bibitem{tr3hironaka}
Yumiko Hironaka.
Spherical functions and local densities on hermitian forms.
\href{https://doi.org/10.2969/jmsj/05130553}
{\emph{Journal of the Mathematical Society of Japan} 51 (1999), 553--581}.
\href{https://www.jstage.jst.go.jp/article/jmath1948/51/3/51_3_553/_pdf/-char/ja\#page=21}
{Formula (3.1), p.~572}, and
\href{https://www.jstage.jst.go.jp/article/jmath1948/51/3/51_3_553/_pdf/-char/ja\#page=25}
{Corollary~3.3(i), p.~576}.

\bibitem{tr3waldron}
Shayne Waldron.
Generalized Welch bound equality sequences are tight frames.
\href{https://doi.org/10.1109/TIT.2003.815788}
{\emph{IEEE Transactions on Information Theory} 49 (2003), 2307--2309}.
Author's version,
\href{https://www.math.auckland.ac.nz/~waldron/Preprints/Welch-bound/CLN-02-481.pdf\#page=2}
{Theorem~1.3} and
\href{https://www.math.auckland.ac.nz/~waldron/Preprints/Welch-bound/CLN-02-481.pdf\#page=3}
{\S2}.

\bibitem{rk3gaudron}
\'Eric Gaudron.
Pentes des fibr\`es vectoriels ad\'eliques sur un corps global.
\emph{Rendiconti del Seminario Matematico della Universit\`a di Padova}
119 (2008), 21--95.
arXiv v2, \href{https://arxiv.org/pdf/math/0605408v2\#page=27}{Proposition~4.22},
\href{https://arxiv.org/pdf/math/0605408v2\#page=34}{Proposition~5.13} and
\href{https://arxiv.org/pdf/math/0605408v2\#page=38}{Definition~6.1}.

\bibitem{tr3chenHN}
Huayi Chen.
Convergence des polygones de Harder--Narasimhan.
\href{https://doi.org/10.24033/msmf.432}
{\emph{M\'emoires de la Soci\'et\'e Math\'ematique de France} (N.S.) 120 (2010)}.
\href{https://www.numdam.org/item/10.24033/msmf.432.pdf\#page=36}
{Equation~(20), p.~34}, and
\href{https://www.numdam.org/item/10.24033/msmf.432.pdf\#page=44}
{Proposition~2.2.4, p.~42}.

\bibitem{lt4bostchen}
Jean-Beno\^it Bost and Huayi Chen.
Concerning the semistability of tensor products in Arakelov geometry.
\emph{Journal de Math\'ematiques Pures et Appliqu\'ees} 99 (2013), 436--488.
\href{https://webusers.imj-prg.fr/~huayi.chen/Recherche/bost_chen.pdf\#page=3}
{\S\S1.2--1.3, the Harder--Narasimhan polygon}, and
\S\S2.2--2.4, especially
\href{https://webusers.imj-prg.fr/~huayi.chen/Recherche/bost_chen.pdf\#page=22}
{(36), p.~457},
\href{https://webusers.imj-prg.fr/~huayi.chen/Recherche/bost_chen.pdf\#page=23}
{Definition~2.3, p.~458}, and
\href{https://webusers.imj-prg.fr/~huayi.chen/Recherche/bost_chen.pdf\#page=26}
{\S2.6, p.~461}.

\bibitem{depth-banaszczyk}
Wojciech Banaszczyk.
New bounds in some transference theorems in the geometry of numbers.
\href{https://doi.org/10.1007/BF01445125}
{\emph{Mathematische Annalen} 296 (1993), 625--635, Theorem~(2.2)}.

\bibitem{lt4bost-jets}
Jean-Beno\^it Bost.
Algebraic leaves of algebraic foliations over number fields.
\href{https://doi.org/10.1007/s10240-001-8191-3}
{\emph{Publications Math\'ematiques de l'IH\'ES} 93 (2001), 161--221}.
\href{https://www.numdam.org/item/PMIHES_2001__93__161_0.pdf\#page=32}
{Proposition~3.9(2), p.~192}, and
\href{https://www.numdam.org/item/PMIHES_2001__93__161_0.pdf\#page=39}
{Propositions~4.5--4.6, pp.~199--201}.

\bibitem{lt4krattenthaler-muller}
Christian Krattenthaler and Thomas W. M\"uller.
Truncated versions of Dwork's lemma for exponentials of power series
and \(p\)-divisibility of arithmetic functions.
\href{https://doi.org/10.1016/j.aim.2015.05.018}
{\emph{Advances in Mathematics} 283 (2015), 489--529}.
\href{https://arxiv.org/pdf/1412.7014v2\#page=4}
{arXiv v2, Theorem~2, (2.4)--(2.6), pp.~4--5}.

\bibitem{lt4david-philippon}
Sinnou David and Patrice Philippon.
Minorations des hauteurs normalis\'ees des sous-vari\'et\'es des tores.
\href{https://www.numdam.org/item/ASNSP_1999_4_28_3_489_0/}
{\emph{Annali della Scuola Normale Superiore di Pisa, Classe di Scienze}
(4) 28 (1999), 489--543}.
\href{https://www.numdam.org/item/ASNSP_1999_4_28_3_489_0.pdf\#page=40}
{Proposition~4.10, p.~527}.

\bibitem{lt4chen-finite-hs}
Huayi Chen.
Explicit uniform estimation of rational points I. Estimation of heights.
\href{https://doi.org/10.1515/CRELLE.2011.138}
{\emph{Journal f\"ur die reine und angewandte Mathematik} 668 (2012),
59--88}.
\href{https://webusers.imj-prg.fr/~huayi.chen/Recherche/effunifI.pdf\#page=22}
{Ahead-of-print version, Theorem~4.8, p.~22}.

\bibitem{lt4cayley}
Sergio Caracciolo, Alan D. Sokal and Andrea Sportiello.
Algebraic/combinatorial proofs of Cayley-type identities for derivatives
of determinants and pfaffians.
\href{https://doi.org/10.1016/j.aam.2012.12.001}
{\emph{Advances in Applied Mathematics} 50 (2013), 474--594}.
\href{https://discovery.ucl.ac.uk/id/eprint/10101899/1/Sokal_1-s2.0-S0196885812001315-main.pdf\#page=4}
{Theorem~2.1, (2.1), p.~477}.

\bibitem{lt4esnault-viehweg}
H\'el\`ene Esnault and Eckart Viehweg.
\emph{Lectures on Vanishing Theorems}.
DMV Seminar 20, Birkh\"auser, 1992.
\href{https://doi.org/10.1007/978-3-0348-8600-0}{Publisher's edition}.
\href{https://page.mi.fu-berlin.de/esnault/books/esvibuch.pdf\#page=50}
{Lemma~5.8(a), p.~47}, and
\href{https://page.mi.fu-berlin.de/esnault/books/esvibuch.pdf\#page=52}
{Corollary~5.12(c), p.~49}.

\bibitem{lt4demailly}
Jean-Pierre Demailly.
Extension of holomorphic functions defined on non reduced analytic
subvarieties.
In \emph{The Legacy of Bernhard Riemann after One Hundred and Fifty Years},
vol.~I, Advanced Lectures in Mathematics 35, International Press, 2016,
pp.~191--222.
\href{https://arxiv.org/pdf/1510.05230v1\#page=3}
{arXiv v1, (2.4)--(2.5), (2.7)--(2.8) and Remark~(2.9)(a), pp.~3--5}.

\bibitem{ar4-harish-chandra}
Harish-Chandra.
Discrete series for semisimple Lie groups. II. Explicit determination of the characters.
\href{https://doi.org/10.1007/BF02392813}{\emph{Acta Mathematica} 116 (1966), 1--111}.
\href{https://repository.ias.ac.in/30833/1/336.pdf\#page=96}
{\S\,41, Theorem~16, p.~96}.

\bibitem{ar4-atiyah-schmid}
Michael F. Atiyah and Wilfried Schmid.
A geometric construction of the discrete series for semisimple Lie groups.
\emph{Inventiones Mathematicae} 42 (1977), 1--62.
\href{https://www.math.utah.edu/~ptrapa/math-library/schmid/atiyah-schmid.pdf}
{Sections~3 and~8, especially (3.10), (3.12) and Theorem~(8.1)}.

\bibitem{ar4-shimura-mass}
Wee Teck Gan, Jonathan P. Hanke and Jiu-Kang Yu.
On an exact mass formula of Shimura.
\href{https://doi.org/10.1215/S0012-7094-01-10716-3}{\emph{Duke Mathematical Journal} 107 (2001), 101--131}.
\href{https://web.math.princeton.edu/~jonhanke/Web-02/Mass-Formula-of-Shimura/shimura-mass.pdf}
{Proposition~2.13 and Section~4, especially (4.2)--(4.3) and Propositions~4.4--4.5}.

\bibitem{bk-local}
Spencer Bloch and Kazuya Kato.
\(L\)-functions and Tamagawa numbers of motives.
In \emph{The Grothendieck Festschrift}, vol.~I,
Progress in Mathematics 86, Birkh\"auser, 1990, 333--400.
\href{https://math.stanford.edu/~conrad/BSDseminar/refs/BKTamagawa.pdf\#page=23}
{Lemma~3.8.1, Corollary~3.8.4 and Example~3.9}.

\bibitem{berger-wach}
Laurent Berger.
Limites de repr\'esentations cristallines.
\href{https://doi.org/10.1112/S0010437X04000879}
{\emph{Compositio Mathematica} 140 (2004), 1473--1498}.
\href{https://perso.ens-lyon.fr/laurent.berger/articles/article06.pdf}
{Sections~II.1--II.3, III.2--III.4 and IV.1}.

\bibitem{chang-diamond}
Seunghwan Chang and Fred Diamond.
Extensions of rank one \((\varphi,\Gamma)\)-modules and crystalline
representations. \href{https://doi.org/10.1112/S0010437X1000504X}
{\emph{Compositio Mathematica} 147 (2011), 375--427}.
\href{https://arxiv.org/pdf/0910.2371v2\#page=41}
{Example~7.1 and Lemma~7.2}.

\bibitem{zhu-labelled}
Hui June Zhu.
Crystalline representations of \(G_{\mathbf Q_{p^a}}\) with coefficients.
\href{https://arxiv.org/pdf/0807.1078v3\#page=3}
{arXiv:0807.1078v3, Theorems~B and~4.15}.

\bibitem{cgm-integral}
Pavel Čoupek, Evangelia Gazaki and Adriano Marmora.
Exactness property of Breuil--Kisin functors and Bloch--Kato Selmer
groups. \href{https://arxiv.org/pdf/2603.26035v3\#page=24}
{Proposition~3.4} and \href{https://arxiv.org/pdf/2603.26035v3\#page=28}
{arXiv:2603.26035v3 (21 October 2026), Theorem~3.12(2) and
Corollary~3.16}.

\bibitem{lt4barquero}
Adrian Barquero-Sanchez, Lindsay Cadwallader, Olivia Cannon,
Tyler Genao and Riad Masri.
Faltings heights of CM elliptic curves and special Gamma values.
\href{https://doi.org/10.1007/s40993-017-0077-7}
{\emph{Research in Number Theory} 3 (2017), article~13}.
Theorem~1.3, (5.13), and (6.2).

\bibitem{siegel-kron}
Carl Ludwig Siegel.
\emph{On Advanced Analytic Number Theory}.
Tata Institute of Fundamental Research, Lectures on Mathematics 23,
original notes 1961, online reissue.
\href{https://mathweb.tifr.res.in/Documents/Publications/Lectures/tifr23.pdf\#page=17}
{Chapter~I, \S1, Theorem~1, p.~13}.

\bibitem{deshalit-artin}
Ehud de Shalit.
Artin \(L\) Functions.
In \emph{An Introduction to the Langlands Program}, edited by
Joseph Bernstein and Stephen Gelbart, Birkh\"auser Boston, 2004,
pp.~73--87.
\href{https://doi.org/10.1007/978-0-8176-8226-2_4}
{\S\S3.2, 3.4 and 4.2}.

\bibitem{rosser-schoenfeld}
J. Barkley Rosser and Lowell Schoenfeld.
Approximate formulas for some functions of prime numbers.
\href{https://doi.org/10.1215/ijm/1255631807}
{\emph{Illinois Journal of Mathematics} 6 (1962), 64--94}.
\href{https://denisevellachemla.eu/Rosser-Schoenfeld-1962.pdf\#page=8}
{Theorems~9 and~12, (3.32) and (3.35), p.~71}.

\bibitem{dlmf}
Tom M. Apostol.
Zeta and related functions. In \emph{NIST Digital Library of
Mathematical Functions}, Chapter~25.
For zeta see the \href{https://dlmf.nist.gov/25.2.i}{continuation},
\href{https://dlmf.nist.gov/25.2.E3}{alternating series},
\href{https://dlmf.nist.gov/25.2.E4}{Laurent expansion},
\href{https://dlmf.nist.gov/25.2.E11}{Euler product},
\href{https://dlmf.nist.gov/25.4.E2}{reflection formula},
\href{https://dlmf.nist.gov/25.6.E1}{value at zero}
and \href{https://dlmf.nist.gov/25.10.i}{zero distribution}.
For Dirichlet \(L\)-functions see the
\href{https://dlmf.nist.gov/25.15.E1}{series and continuation},
\href{https://dlmf.nist.gov/25.15.E2}{Euler product},
\href{https://dlmf.nist.gov/25.15.E4}{change of conductor},
\href{https://dlmf.nist.gov/25.15.E5}{functional equation}
and \href{https://dlmf.nist.gov/25.15.E9}{nonvanishing at one}.
For the gamma function see R.\ A.\ Askey and R.\ Roy, Chapter~5,
\href{https://dlmf.nist.gov/5.2.i}{\S5.2(i)}.

\bibitem{z6montgomery-vaughan}
Hugh L. Montgomery and Robert C. Vaughan.
\emph{Multiplicative Number Theory I. Classical Theory}.
Cambridge Studies in Advanced Mathematics 97. Cambridge University Press, 2007.
\href{https://personal.science.psu.edu/rcv4/personal/Publications/MNTI/06.0_pp_35_75_The_elementary_theory_of_arithmetic_functions.pdf\#page=16}
{Theorem~2.7(a), p.~50},
\href{https://personal.science.psu.edu/rcv4/personal/Publications/MNTI/08.0_pp_108_136_Primes_in_arithmetic_progressions_I.pdf\#page=9}
{Theorem~4.4, p.~116},
\href{https://personal.science.psu.edu/rcv4/personal/Publications/MNTI/14.0_pp_326_357_Analytic_properties_of_the_zeta_function_and_L-functions.pdf\#page=7}
{(10.17) and Corollary~10.8, pp.~332--333}, and
\href{https://personal.science.psu.edu/rcv4/personal/Publications/MNTI/15.0_pp_358_396_Primes_in_arithmetic_progressions_II.pdf\#page=3}
{Theorem~11.3, p.~360}.

\bibitem{z6kaczorowski-pintz}
J. Kaczorowski and J. Pintz.
Oscillatory properties of arithmetical functions. I.
\href{https://doi.org/10.1007/BF01949062}
{\emph{Acta Mathematica Hungarica} 48 (1986), 173--185}.
\href{https://personal.math.ubc.ca/~gerg/teaching/592-Fall2018/papers/1986.Kaczorowski_1.pdf\#page=3}
{Theorem~1 and (2.2), p.~175}.

\bibitem{z6kaczorowski-pintz-ii}
J. Kaczorowski and J. Pintz.
Oscillatory properties of arithmetical functions. II.
\href{https://doi.org/10.1007/BF01951008}
{\emph{Acta Mathematica Hungarica} 49 (1987), 441--453}.
\href{https://personal.math.ubc.ca/~gerg/teaching/592-Fall2018/papers/1987.Kaczorowski_2.pdf\#page=1}
{Theorem, (1.5) and (1.7), pp.~441--442}.

\bibitem{real-branch-knapp-vogan}
Anthony W. Knapp and David A. Vogan, Jr.
\emph{Cohomological Induction and Unitary Representations}.
Princeton Mathematical Series 45, Princeton University Press, 1995.
\href{https://www.math.stonybrook.edu/~aknapp/books/kv.html}{Book page}.

\bibitem{real-branch-prudhon}
Nicolas Prudhon.
Translation of Dolbeault representations.
\href{https://www.heldermann.de/JLT/JLT29/JLT294/jlt29050.htm}{\emph{Journal of Lie Theory} 29 (2019), 1103--1118}.
\href{https://prudhon.perso.math.cnrs.fr/fichiers/GeomTranslPrinciple-dolbeault.pdf}{Proposition~3.6 and (14)}.

\bibitem{real-branch-bernstein-kroetz}
Joseph Bernstein and Bernhard Kr\"otz.
Smooth Fr\'echet globalizations of Harish-Chandra modules.
\emph{Israel Journal of Mathematics} 199 (2014), 45--111.
\href{https://arxiv.org/pdf/0812.1684}{Theorems~1.1 and 10.6}.

\bibitem{real-branch-mno}
Alexander Molev, Maxim Nazarov and Grigori Olshanski.
Yangians and classical Lie algebras.
\emph{Russian Mathematical Surveys} 51 (1996), 205--282.
\href{https://arxiv.org/pdf/hep-th/9409025}{Proposition~2.7(4) and Remark~2.11}.

\bibitem{real-branch-sun}
Binyong Sun.
Matrix coefficients of cohomologically induced representations.
\href{https://doi.org/10.1112/S0010437X06002508}{\emph{Compositio Mathematica} 143 (2007), 201--221}.
See Theorem~6.3 and (25)--(26).

\bibitem{p2bernstein-unitary-induction}
Joseph N. Bernstein.
P-invariant distributions on \(\mathrm{GL}(N)\) and the classification
of unitary representations of \(\mathrm{GL}(N)\) (non-Archimedean case).
In \emph{Lie Group Representations II}, Lecture Notes in Mathematics
1041, Springer, 1984, 50--102.
\href{https://www.math.tau.ac.il/~bernstei/Publication_list/publication_texts/Bernstein-P-invar-SLN.pdf}{Author's copy}.

\bibitem{p2zelevinsky-induction-ring}
Andrei V. Zelevinsky.
Induced representations of reductive \(p\)-adic groups. II.
On irreducible representations of \(\mathrm{GL}(n)\).
\emph{Annales Scientifiques de l'\'Ecole Normale Sup\'erieure}
13 (1980), 165--210.
\href{https://numdam.org/item/10.24033/asens.1379.pdf}{Published article}.

\bibitem{p2platonov-strong-approximation}
Vladimir Platonov and Andrei Rapinchuk.
\emph{Algebraic Groups and Number Theory}.
Academic Press, 1994.
\href{https://uva.theopenscholar.com/files/andrei-rapinchuk/files/agnt_english.pdf}{Author's copy}.

\bibitem{p2wu-unitary-tower}
Chenyan Wu.
Irreducibility of theta lifting for unitary groups.
\emph{Journal of Number Theory} 133 (2013), 3296--3318.
\href{https://arxiv.org/pdf/1409.0770\#page=9}{Author's version}.

\bibitem{p2yamana-singular-siegel-weil}
Shunsuke Yamana.
On the Siegel--Weil formula: the case of singular forms.
\emph{Compositio Mathematica} 147 (2011), 1003--1021.
\href{https://doi.org/10.1112/S0010437X11005379}{Published article}.

\bibitem{p2bernstein-lapid-eisenstein}
Joseph Bernstein and Erez Lapid.
On the meromorphic continuation of Eisenstein series.
\emph{Journal of the American Mathematical Society} 37 (2024), 187--234.
\href{https://www.math.tau.ac.il/~bernstei/Publication_list/publication_texts/2023-10-15-Eisenstein-series-JAMS.pdf}{Author's copy}.

\bibitem{p2miller-schmid-decay}
Stephen D. Miller and Wilfried Schmid.
On the rapid decay of cuspidal automorphic forms.
\emph{Advances in Mathematics} 231 (2012), 940--964.
\href{https://arxiv.org/pdf/1106.2149\#page=3}{Theorem~A},
\href{https://arxiv.org/pdf/1106.2149\#page=11}{Theorems~2.23--2.24}.

\bibitem{ar4-graftieaux}
Philippe Graftieaux.
Formal groups and the isogeny theorem.
\href{https://doi.org/10.1215/S0012-7094-01-10614-5}{\emph{Duke Mathematical Journal} 106 (2001), 81--121}.
See Proposition~2.8, Theorem~2.10 and Propositions~2.11 and~2.13
for filtered slopes, cubist models and jets along formal subgroups.

\bibitem{mumford-theta-isogeny}
David Mumford.
On the equations defining abelian varieties. I.
\href{https://doi.org/10.1007/BF01389737}{\emph{Inventiones Mathematicae} 1 (1966), 287--354}.
\href{https://pazuki.perso.math.cnrs.fr/index_fichiers/Mumford66.pdf\#page=16}{Section~1, Theorem~4, pp.~302--303}.

\bibitem{tame-selmer-mazur-rubin}
Barry Mazur and Karl Rubin, with an appendix by Michael Larsen.
\emph{Diophantine stability}.
\href{https://arxiv.org/pdf/1503.04642v3}{Proposition~9.17, Theorem~9.21 and Remark~10.4}.

\bibitem{tame-selmer-smith}
Alexander Smith.
\emph{The distribution of \(\ell^\infty\)-Selmer groups in degree
\(\ell\) twist families II}.
\href{https://arxiv.org/pdf/2207.05143v2\#page=10}{Theorem~2.14}.

\bibitem{paul-trapa-characters}
Annegret Paul and Peter E. Trapa.
One-dimensional representations of \(U(p,q)\) and the Howe correspondence.
\href{https://doi.org/10.1006/jfan.2002.3974}{\emph{Journal of
Functional Analysis} 195 (2002), 129--166}.
\href{https://www.math.utah.edu/~ptrapa/char6_4.pdf\#page=9}{Theorem~4.1}.

\bibitem{bao-sun-smooth-theta}
Yixin Bao and Binyong Sun.
\emph{Coincidence of algebraic and smooth theta correspondences}.
\href{https://arxiv.org/pdf/1611.06298\#page=14}{Proposition~5.1, unitary dual pairs}.

\bibitem{p2paul-real-theta}
Annegret Paul.
First occurrence for the dual pairs \((U(p,q),U(r,s))\).
\href{https://doi.org/10.4153/CJM-1999-029-6}{\emph{Canadian Journal
of Mathematics} 51 (1999), 636--657}.
See Theorem~2.7 and Proposition~3.4.

\bibitem{p2milicic-integrability}
Dragan Mili\v{c}i\'c.
Asymptotic behaviour of matrix coefficients of the discrete series.
\emph{Duke Mathematical Journal} 44 (1977), 59--88.
\href{https://www.math.utah.edu/~ptrapa/math-library/milicic/milicic-asymptotics-duke-1977.pdf\#page=26}{Part~III, Theorem~1 and its corollary}.

\bibitem{p2flip-howe-real}
Roger Howe.
Transcending classical invariant theory.
\href{https://doi.org/10.1090/S0894-0347-1989-0985172-6}{\emph{Journal
of the American Mathematical Society} 2 (1989), 535--552}.
\href{https://www.math.utah.edu/~ptrapa/math-library/howe/howe-transcending-1989.pdf}{Theorems~1 and~2.1}.

\bibitem{p2flip-gan-takeda}
Wee Teck Gan and Shuichiro Takeda.
A proof of the Howe duality conjecture.
\href{https://doi.org/10.1090/jams/839}{\emph{Journal of the American
Mathematical Society} 29 (2016), 473--493}.
\href{https://arxiv.org/pdf/1407.1995v4}{Theorem~1.2}.

\bibitem{p2casselman-spherical}
W. Casselman.
The unramified principal series of \(p\)-adic groups. I. The spherical
function. \emph{Compositio Mathematica} 40 (1980), 387--406.
\href{https://www.numdam.org/article/CM_1980__40_3_387_0.pdf\#page=18}{Proposition~4.1 and Theorem~4.2 (Macdonald)}.

\bibitem{flag-hormann}
Fritz H\"ormann.
Generalized automorphic sheaves and the proportionality principle of
Hirzebruch--Mumford.
\href{https://arxiv.org/pdf/1603.04411v2\#page=32}{Theorem~4.2.1}.

\bibitem{flag-bruinier-howard}
Jan Hendrik Bruinier and Benjamin Howard.
Arithmetic volumes of unitary Shimura varieties.
\href{https://arxiv.org/pdf/2105.11274v3\#page=117}{Theorem~8.6.1}.

\bibitem{flag-guo}
Ziqi Guo.
Modular heights of unitary Shimura varieties III: Proof of the main
theorem. Preprint (2026).
\href{https://arxiv.org/pdf/2606.19158v4\#page=4}{Section~1.1 and Theorem~1.1}.

\bibitem{flag-maillot-rossler}
Vincent Maillot and Damian R\"ossler.
Conjectures on the logarithmic derivatives of Artin \(L\)-functions II.
\href{https://arxiv.org/pdf/1808.03068v1\#page=41}{Theorem~7.26}.

\bibitem{rank-labesse}
Jean-Pierre Labesse.
Changement de base CM et s\'eries discr\`etes.
\href{https://webusers.imj-prg.fr/~michael.harris/CBCM.pdf\#page=42}
{Condition \((*)\) and Corollary~5.3}.

\bibitem{rank-salamanca}
Susana Salamanca--Riba.
On the unitary dual of some classical Lie groups.
\emph{Compositio Mathematica} 68 (1988), 251--303.
\href{https://www.numdam.org/item/CM_1988__68_3_251_0.pdf\#page=4}
{Theorem~1.2 and Proposition~4.3}.

\bibitem{rank-moeglin}
Colette M\oe glin and Jean-Loup Waldspurger.
Le spectre r\'esiduel de \(\mathrm{GL}(n)\).
\emph{Annales scientifiques de l'\'Ecole Normale Sup\'erieure}
22 (1989), 605--674.
\href{https://www.numdam.org/item/ASENS_1989_4_22_4_605_0.pdf\#page=3}
{Theorem of the introduction}.

\bibitem{rank-jacquet-shalika}
Herv\'e Jacquet and Joseph A. Shalika.
On Euler products and the classification of automorphic forms II.
\href{https://doi.org/10.2307/2374050}
{\emph{American Journal of Mathematics} 103 (1981), 777--815}.
\href{https://www.math.columbia.edu/~hj/On\%20Euler\%20products\%20II.pdf\#page=34}
{Theorem~4.4, p.~809}.

\bibitem{rank-kmsw}
Tasho Kaletha, Alberto M\'inguez, Sug Woo Shin and Paul-James White.
Endoscopic classification of representations: inner forms of unitary groups.
\href{https://public.websites.umich.edu/~kaletha/unitary.pdf\#page=150}
{Corollary~3.3.2} and
\href{https://public.websites.umich.edu/~kaletha/unitary.pdf\#page=89}
{Theorem*~1.7.1 with its proved generic, pure-inner scope}.

\bibitem{rank-harris-spherical}
Michael Harris.
Cohomological automorphic forms on unitary groups, II:
period relations and values of \(L\)-functions.
\href{https://webusers.imj-prg.fr/~michael.harris/periodrelations.pdf\#page=25}
{Corollary~2.2.7 and Appendix~2.A, especially Corollaries~2.A.3 and 2.A.11}.

\bibitem{weyl-andersen}
Henning Haahr Andersen.
Filtrations of cohomology modules for Chevalley groups.
\emph{Annales scientifiques de l'\'Ecole Normale Sup\'erieure}
16 (1983), 495--528.
\href{https://www.numdam.org/item/10.24033/asens.1458.pdf\#page=16}
{Propositions~3.2 and~4.1, pp.~505--506 and 509--510}.

\bibitem{weyl-tange}
Rudolf Tange.
A bideterminant basis for a reductive monoid.
\emph{Journal of Pure and Applied Algebra} 216 (2012), 1207--1221.
\href{https://rtange.github.io/research/tableaux.pdf\#page=4}
{Theorem~1.1 and its corollary}.

\bibitem{weyl-lin}
Zongzhu Lin.
Representations of Chevalley groups arising from admissible lattices.
\href{https://www.math.ksu.edu/~zlin/research/adm.pdf\#page=3}
{Sections~1.1 and 2.2--2.3}.

\bibitem{bgk-weight-elimination}
Bhargav Bhatt, Toby Gee and Mark Kisin.
Reduction modulo \(p\) of crystalline Galois representations via
\(\mu_p\)-equivariance. Preprint, arXiv:2607.08660v2 (2026).
\href{https://arxiv.org/pdf/2607.08660v2\#page=77}
{Theorem~5.6.2 for unramified base fields}.

\bibitem{nakahama-norms}
Ryosuke Nakahama.
Norm computation and analytic continuation of vector valued
holomorphic discrete series representations.
\href{https://doi.org/10.5802/jolt.919}{\emph{Journal of Lie Theory}
26 (2016), 927--990}.
\href{https://arxiv.org/pdf/1506.05919\#page=5}
{Theorem~1.2 (Theorems~4.3 and~5.1), the ratio to the Fischer norm}.

\bibitem{content-partitions}
Chi-Fang Chen, Jorge Garza-Vargas and Ramon van Handel.
A new approach to strong convergence II. The classical ensembles.
\href{https://web.math.princeton.edu/~rvan/strongii241216.pdf\#page=29}
{Proof of Lemma~7.4, the multiplicity of a fixed content in a partition}
(2024 preprint).

\bibitem{depth-katz}
Nicholas M. Katz.
Formal groups and \(p\)-adic interpolation.
\emph{Ast\'erisque} 41--42 (1977), 55--65.
\href{https://www.numdam.org/article/AST_1977__41-42__55_0.pdf\#page=9}
{Section~V, height two}.

\bibitem{depth-lubin-tate}
Peter Schneider and Jeremy Teitelbaum.
\(p\)-adic Fourier theory.
\href{https://arxiv.org/pdf/math/0102012\#page=33}
{Appendix, Theorem~(c), the valuation of the fundamental period}.

\bibitem{orbit-kollar}
J\'anos Koll\'ar.
Singularities of pairs.
\href{https://arxiv.org/pdf/alg-geom/9601026v2\#page=35}
{Theorem~6.4, an isolated non-klt point of a divisor in a prescribed numerical class}.

\bibitem{orbit-nadel}
Robert Lazarsfeld.
A short course on multiplier ideals.
\href{https://www.math.stonybrook.edu/~roblaz/Reprints/Lazarsfeld.PCMI.Short.Course.MI.pdf\#page=11}
{Theorem~2.4 and its proof, Nadel vanishing for a nef and big difference}.

\bibitem{affinoid-residue}
Siegfried Bosch.
\emph{Lectures on Formal and Rigid Geometry}.
\href{https://www.math.purdue.edu/~tongliu/seminar/rigid/Bosch.pdf}
{Section~2.2, Corollaries~12--13, residue fields and geometric points of the Tate algebra}.

\bibitem{galois-raec}
Thomas Barnet-Lamb, Toby Gee, David Geraghty and Richard Taylor.
Potential automorphy and change of weight.
\emph{Annals of Mathematics} 179 (2014), 501--609.
\href{https://www.ma.imperial.ac.uk/~tsg/Index_files/potentialautomorphy.pdf\#page=33}
{Theorem~2.1.1(2)--(4), purity away from \(\ell\), labelled Hodge--Tate weights and the Iwahori case at \(\ell\)}.

\bibitem{galois-wd-at-ell}
Thomas Barnet-Lamb, Toby Gee, David Geraghty and Richard Taylor.
Local-global compatibility for \(l=p\), II.
\emph{Annales scientifiques de l'\'{E}cole Normale Sup\'erieure} 47 (2014), 165--179.
\href{https://www.ma.imperial.ac.uk/~tsg/Index_files/localglobalII.pdf\#page=6}
{Theorem~1.1, comparison of the monodromy partitions above \(\ell\)}.

\bibitem{bost-theta-key}
Jean-Beno\^it Bost.
Intrinsic heights of stable varieties and abelian varieties.
\href{https://doi.org/10.1215/S0012-7094-96-08202-2}
{\emph{Duke Mathematical Journal} 82 (1996), 21--70},
Theorem~4.10(v).

\bibitem{skoda-division}
Henri Skoda.
Application des techniques \(L^2\) à la théorie des idéaux
d'une algèbre de fonctions holomorphes avec poids.
\emph{Annales scientifiques de l'École Normale Supérieure}
5 (1972), 545--579.
\href{https://www.numdam.org/article/ASENS_1972_4_5_4_545_0.pdf\#page=11}
{Theorem~1, pp.~555--556, equations~(3.1)--(3.2)}.

\bibitem{flag-hormander-l2}
Jean-Pierre Demailly.
\emph{Analytic methods in algebraic geometry}, Grenoble lectures,
December 2009.
\href{https://www-fourier.ujf-grenoble.fr/~demailly/manuscripts/analmeth.pdf}
{\S5.A, Theorem~5.1}, the complete Kähler estimate with the inverse
curvature operator on bundle-valued forms.

\bibitem{atobe-harder-crystalline}
Hiraku Atobe, Masataka Chida, Tomoyoshi Ibukiyama, Hidenori Katsurada
and Takuya Yamauchi. Harder's conjecture II.
\href{https://doi.org/10.1016/j.aim.2026.111233}
{\emph{Advances in Mathematics} 503 (2026), article 111233}.
\href{https://arxiv.org/pdf/2306.07582v4\#page=25}
{Extended version, Lemma~7.6}, on subquotient extensions over the
residue field of a short crystalline lattice.

\bibitem{lan-arc-theta}
Kai-Wen Lan.
\emph{Arithmetic Compactifications of PEL-Type Shimura Varieties}.
Princeton University Press, 2013.
\href{https://www.kwlan.org/articles/cpt-PEL-type-thesis-revision-dense.pdf\#page=96}
{Equation~(4.3.1.7), the Fourier coefficient recurrence},
\href{https://www.kwlan.org/articles/cpt-PEL-type-thesis-revision-dense.pdf\#page=105}
{Theorem~4.4.16, Mumford's equivalence with the polarized degeneration data},
and \href{https://www.kwlan.org/articles/cpt-PEL-type-thesis-revision-dense.pdf\#page=116}
{\S4.5.3, torsion over the abelian locus}.

\bibitem{pasten-cm-norm}
Hector Pasten.
\emph{Shimura curves and the abc conjecture}.
\href{https://arxiv.org/pdf/1705.09251v4\#page=58}
{\S\S13--14, Theorem~14.1 and Proposition~14.2},
on Heegner point heights and norms of integral differentials.

\bibitem{richter-westerholt-sturm}
Olav K. Richter and Martin Westerholt-Raum.
Sturm bounds for Siegel modular forms.
\href{https://doi.org/10.1007/s40993-015-0008-4}
{\emph{Research in Number Theory} 1 (2015), article~5}.
\href{https://arxiv.org/pdf/1501.07733v1\#page=5}
{Proof of Proposition~2.4}, torsion specialization of Jacobi forms.

\bibitem{de-jong-shokrieh}
Robin de Jong and Farbod Shokrieh.
Faltings height and N\'eron--Tate height of a theta divisor.
\href{https://doi.org/10.1112/S0010437X21007661}
{\emph{Compositio Mathematica} 158 (2022), 1--32}.
Equations~(12.1)--(12.2) give the Faltings metric and the
probability \(L^2\)-metric. Theorem~13.1 and Remark~13.1 give
Bost's key formula for a principally polarized Moret-Bailly model
in these conventions.

\bibitem{cartier-display}
Thomas Zink.
The display of a formal \(p\)-divisible group.
\emph{Astérisque} 278 (2002), 127--248.
\href{https://www.numdam.org/article/AST_2002__278__127_0.pdf\#page=113}
{\S4.3, equations~(216)--(218), the biextension associated with a pairing of displays}.

\bibitem{i2-neukirch}
J\"urgen Neukirch.
\emph{Algebraic Number Theory}.
\href{https://doi.org/10.1007/978-3-662-03983-0}{Grundlehren der mathematischen
Wissenschaften 322, Springer, Berlin, 1999}.

\bibitem{i2-washington}
Lawrence C. Washington.
\emph{Introduction to Cyclotomic Fields}, 2nd ed.
\href{https://doi.org/10.1007/978-1-4612-1934-7}{Graduate Texts in Mathematics 83,
Springer, New York, 1997}.

\bibitem{f1-blggt-pacw}
Thomas Barnet-Lamb, Toby Gee, David Geraghty and Richard Taylor.
Potential automorphy and change of weight.
\href{https://doi.org/10.4007/annals.2014.179.2.3}
{\emph{Annals of Mathematics} 179 (2014), 501--609}.
\href{https://arxiv.org/pdf/1010.2561v3\#page=33}
{arXiv v3, Theorem~2.1.1, pp.~33--34}.

\bibitem{f1-caraiani-lp}
Ana Caraiani.
Monodromy and local-global compatibility for \(l=p\).
\href{https://doi.org/10.2140/ant.2014.8.1597}
{\emph{Algebra \& Number Theory} 8 (2014), 1597--1646}.
\href{https://arxiv.org/pdf/1202.4683v1\#page=1}
{arXiv v1, Theorem~1.1, p.~1}.

\bibitem{s1-thorner-zaman-chebotarev}
Jesse Thorner and Asif Zaman.
A unified and improved Chebotarev density theorem.
\href{https://doi.org/10.2140/ant.2019.13.1039}{\emph{Algebra \& Number
Theory} 13 (2019), no.~5, 1039--1068}.
\href{https://arxiv.org/pdf/1803.02823v3\#page=1}{arXiv v3, (1.4)},
\href{https://arxiv.org/pdf/1803.02823v3\#page=3}{(1.11) and Theorem~1.4},
\href{https://arxiv.org/pdf/1803.02823v3\#page=7}{Section~2}.

\bibitem{sr1-gross-minimal}
Benedict H. Gross.
Minimal models for elliptic curves with complex multiplication.
\emph{Compositio Math.} 45 (1982), 155--164.
\href{https://www.numdam.org/item/CM_1982__45_2_155_0.pdf\#page=5}
{\S3, the curve \(A(p)\) before Proposition~3.2: good reduction outside \(p\), minimal discriminant \((-p^3)\)}.

\bibitem{sr1-serre-local}
Jean-Pierre Serre.
\emph{Local Fields}.
\href{https://doi.org/10.1007/978-1-4757-5673-9}{Graduate Texts in Mathematics 67,
Springer, New York, 1979}, Ch.~III, \S6.

\bibitem{sr1-silverman-aec}
Joseph H. Silverman.
\emph{The Arithmetic of Elliptic Curves}, 2nd ed.
\href{https://doi.org/10.1007/978-0-387-09494-6}{Graduate Texts in Mathematics 106,
Springer, New York, 2009}.

\bibitem{sr1-silverman-advanced}
Joseph H. Silverman.
\emph{Advanced Topics in the Arithmetic of Elliptic Curves}.
\href{https://doi.org/10.1007/978-1-4612-0851-8}{Graduate Texts in Mathematics 151,
Springer, New York, 1994}.

\bibitem{sr1-cox}
David A. Cox.
\emph{Primes of the Form \(x^2+ny^2\)}, 2nd ed.
\href{https://doi.org/10.1002/9781118400722}{John Wiley \& Sons, Hoboken, NJ, 2013}.

\bibitem{s2b-atiyah-macdonald}
Michael F. Atiyah and Ian G. Macdonald.
\emph{Introduction to Commutative Algebra}.
Addison-Wesley, Reading, Mass., 1969.
Ch.~5: Prop.~5.1, Cor.~5.3, Cor.~5.4, Cor.~5.8 and Thm.~5.10.

\bibitem{p7-satake-spherical}
I. Satake.
Theory of spherical functions on reductive algebraic groups over
\(\mathfrak p\)-adic fields.
\emph{Publications Math\'ematiques de l'IH\'ES} 18 (1963), 5--69.
\href{https://www.numdam.org/item/PMIHES_1963__18__5_0.pdf}{Theorem~1, and \S8 for \(\mathrm{GL}(n)\)}.

\bibitem{p7-mvw}
C. M\oe glin, M.-F. Vign\'eras and J.-L. Waldspurger.
\emph{Correspondances de Howe sur un corps \(p\)-adique}.
Lecture Notes in Mathematics 1291, Springer, 1987.
\href{https://doi.org/10.1007/BFb0082712}{Chapter~5, Lemma~I.8, Theorem~I.10 and Proposition~I.11}.

\bibitem{f2-huber-kings}
Annette Huber and Guido Kings.
Bloch--Kato conjecture and main conjecture of Iwasawa theory for
Dirichlet characters.
\emph{Duke Mathematical Journal} 119 (2003), 393--464.
\href{https://arxiv.org/pdf/math/0101071v2\#page=13}
{arXiv v2, Theorem~1.4.1, p.~13}.

\bibitem{f2-ferrero-washington}
Bruce Ferrero and Lawrence C. Washington.
The Iwasawa invariant \(\mu_p\) vanishes for abelian number fields.
\href{https://doi.org/10.2307/1971116}
{\emph{Annals of Mathematics} 109 (1979), 377--395}.

\bibitem{f2-dlmf24}
Karl Dilcher.
Bernoulli and Euler polynomials. In \emph{NIST Digital Library of
Mathematical Functions}, Chapter~24.
See the \href{https://dlmf.nist.gov/24.2.E5}{expansion (24.2.5)} and
the \href{https://dlmf.nist.gov/24.10.E1}{von Staudt--Clausen theorem (24.10.1)}.

\bibitem{p4-lang-algebra}
Serge Lang.
\emph{Algebra}, revised third edition.
Graduate Texts in Mathematics 211, Springer, 2002.
\href{https://doi.org/10.1007/978-1-4613-0041-0}{doi:10.1007/978-1-4613-0041-0}.
Chapter~XVII, \S3 (the density theorem).

\bibitem{s2a-guralnick-lorenz}
Robert M. Guralnick and Martin Lorenz.
Orders of finite groups of matrices.
In \emph{Groups, Rings and Algebras}, Contemporary Mathematics 420,
American Mathematical Society (2006), 141--161.
\href{https://arxiv.org/pdf/math/0511191v1\#page=2}
{arXiv v1, Theorem~1 (Minkowski) and (1), p.~2}.

\end{thebibliography}
\endgroup
